\section{Native rational orders and weighted principal balance on the original curve}

\begin{theorem}[The original rational function's actual native orders]
\label{thm:fc-original-rational-denominator-orders}
\leanok
\lean{EnvelopeProperty.CurveRationalPrincipalBalance.ord_eq_denominator_difference, EnvelopeProperty.CurveRationalPrincipalBalance.ord_eq_normal_denominator_difference}
\uses{thm:fc-rational-denominator-localization, thm:fc-curve-denominator-cartier, thm:fc-curve-cartier-native-local-orders}
Let $X$ be the original integral locally Noetherian scheme and
$0\ne f\in K(X)$ the original rational function. Whenever its two
actual denominator ideals are Cartier,
\[
 \operatorname{ord}_x(f)
 =\operatorname{ord}_x D_X(f^{-1})-\operatorname{ord}_x D_X(f)
\]
for EVERY original point $x$, where all orders are the native original
scheme orders. On an actual normal curve of dimension at most one the
two Cartier facts are derived from the original geometry.
\end{theorem}
\begin{proof}
Take an actual original affine Cartier equation $c$ of $D_X(f)$ near $x$.
The genuine original fraction-field numerator $b=cf$ is an element of
the same original affine ring. Reciprocal multiplication is an actual
linear isomorphism of the two denominator ideals, so their actual
principal generators satisfy $D_X(f^{-1})(U)=(b)$. More explicitly,
surjectivity writes each reciprocal denominator as the image of a multiple
of $c$, and the literal field coefficients give the same multiple of $b$.
The reverse containment is the genuine reciprocal-denominator membership.
Both original equations are nonzero. The established arbitrary-equation
comparison identifies their native orders with the actual Cartier local
orders. Native multiplicativity applied to the original identity $b=cf$
gives the formula, including the zero native order away from codimension
one. The normal-curve Cartier producer supplies the two Cartier facts in
the asserted geometric specialization.
\end{proof}

\begin{definition}[The actual residue-weighted original rational order]
\label{def:fc-original-weighted-rational-order}
\leanok
\lean{EnvelopeProperty.CurveRationalPrincipalBalance.weightedOrder, EnvelopeProperty.CurveRationalPrincipalBalance.hasFiniteSupport_ord, EnvelopeProperty.CurveRationalPrincipalBalance.hasFiniteSupport_weightedOrder}
For an original structure morphism $p:X\to\operatorname{Spec}k$, weight
its native order by the actual original residue-field degree
\[
 w_{p,f}(x)=\operatorname{ord}_x(f)[\kappa(x):k],
\]
using the literal original coefficient map to $\kappa(x)$.
On an original normal Noetherian integral curve of dimension at most one,
both the native rational-order function and this actual residue-weighted
function have finite support, derived from the two true Cartier supports.
\end{definition}

\begin{theorem}[True principal weighted balance on the whole original curve]
\label{thm:fc-original-curve-principal-weighted-balance}
\leanok
\lean{EnvelopeProperty.CurveRationalPrincipalBalance.finsum_weightedOrder_eq_zero}
\uses{thm:fc-original-rational-denominator-orders, def:fc-original-weighted-rational-order, thm:fc-scheme-denominator-reciprocal-iso, thm:fc-curve-cartier-weighted-degree}
For ANY original integral normal projective curve over ANY original field
and ANY nonzero original rational function,
\[
 \sum_{x\in X} \operatorname{ord}_x(f)[\kappa(x):k]=0.
\]
The sum is over EVERY original curve point, with the actual original
residue fields; no separability, perfectness, algebraic closure, finite
support, divisor, degree, Cartier, balancing or order certificate is
assumed.
\end{theorem}
\begin{proof}
The original projective morphism derives properness, finite type, local
Noetherianity, compactness and Noetherianity of the original scheme. The
true denominator Cartier producer gives the two original Cartier ideals.
Their actual original closed supports are finite: their nonzero original
generic germs give proper support, and the genuine curve-support theorem
makes the actual closed subschemes finite. Thus the actual two weighted
Cartier-order functions have finite support on the WHOLE original curve.

Whole original reciprocal multiplication gives an isomorphism of the
same two original denominator modules. Genuine full Euler characteristic
is invariant under this whole original isomorphism. The true original
Cartier-degree formula identifies their normalized Euler characteristics
with minus their actual residue-weighted native-order sums. Therefore
these sums are equal. The preceding pointwise native order formula,
multiplied by the SAME original residue-field degree, identifies the
weighted rational-order function with their difference. Finite-support
sum additivity now gives zero, with every auxiliary finiteness and Cartier
fact derived internally from the stated original geometry.
\end{proof}
